
\Problem{Ordenación por Selección con Iteradores}

\Statement

El algoritmo de \textbf{ordenación por selección} consiste en, de forma iterativa, encontrar el elemento mínimo de la parte no ordenada y colocarlo al principio de esa parte. Concretamente:

\begin{enumerate}
\item Se busca el elemento mínimo del rango \verb|[begin, end)| y se intercambia con el elemento apuntado por \verb|begin|.
\item Seguidamente, se busca el mínimo del rango que va desde el siguiente a \verb|begin| hasta \verb|end|, y se intercambia con el segundo elemento.
\item Se repite hasta que queda un solo elemento por procesar.
\end{enumerate}

Escribe una función con la cabecera
\begin{verbatim}
/**
 * @brief Ordena un rango de una lista usando selección.
 *
 * @param begin Iterador al principio del rango.
 * @param end Iterador al final del rango (no incluido).
 * @pre begin y end delimitan un rango válido de una lista.
 * @post El rango [begin, end) queda ordenado ascendentemente.
 */
void ordenacio_seleccio(list<int>::iterator begin,
                        list<int>::iterator end);
\end{verbatim}
que ordene el rango \verb|[begin, end)| de una lista de enteros usando el algoritmo de ordenación por selección.

Por ejemplo, si la lista contiene \verb|5 3 1 4 2|, después de llamar a la función la lista debe contener \verb|1 2 3 4 5|.

\Input

La entrada consiste en una secuencia de enteros separados por espacios, que se lee hasta el final de la entrada.

\Output

La salida es una línea con los elementos de la lista ordenada ascendentemente, separados por espacios.

\Observation

En los ficheros públicos (icono del gatito) encontrarás: \verb|main.cc| (el programa principal, con la entrada/salida hecha), \verb|ordenacio.cc| (donde debes implementar la función) y un \verb|Makefile|.

Debes implementar \verb|ordenacio_seleccio| en el fichero \verb|ordenacio.cc| y enviar \textbf{únicamente} ese fichero.

\SampleTwoCol
